\documentclass[10pt,a4paper]{amsart}
\usepackage[margin=19mm]{geometry}
\usepackage{amsmath,amssymb,mathtools}
\usepackage[scr=rsfs,cal=euler]{mathalfa}
\usepackage{xfrac}
\usepackage[unicode,hidelinks,bookmarks=false]{hyperref}
\newcommand{\ie}{i.e.\ }
\newcommand{\cf}{cf.\ }
\newcommand{\resp}{respectively\ }
\newcommand{\Subsection}{\subsection}
\providecommand{\nobreakdash}{\nobreak\hbox{-}}
\setlength{\emergencystretch}{2em}
\multlinegap0pt
\usepackage{bm}
\usepackage{bbm}
\usepackage{dsfont}
\usepackage{xspace}
\usepackage{cancel}
\usepackage{mathtools}
\usepackage{amsthm}
\makeatletter
\@ifundefined{theorem}{\newtheorem{theorem}{Theorem}}{}
\@ifundefined{lemma}{\newtheorem{lemma}[theorem]{Lemma}}{}
\makeatother
\DeclareMathOperator{\Res}{Res}
\newcommand{\NN}{\mathbb{N}}
\newcommand{\Q}{\mathbb{Q}}
\renewcommand{\geq}{\geqslant}
\renewcommand{\leq}{\leqslant}
\title[Solubility of a resultant equation and applications]{Solubility of a resultant equation and applications}
\author{Tim Browning, Stephanie Chan}
\date{Source: arXiv:2411.09264v2 (CC BY 4.0)}
\begin{document}
\maketitle

\noindent\emph{Source and licence.} Solubility of a resultant equation and applications. Version arXiv:2411.09264v2 (CC BY 4.0), distributed under the Creative Commons Attribution 4.0 International licence, as recorded in the arXiv licence metadata for this exact version and shown on the abstract page https://arxiv.org/abs/2411.09264v2. Full rights notice: licenses/arxiv-2411.09264-CC-BY-4.0.txt.\par\smallskip

\noindent\textit{Editorial scope.} Counted unit: the Lemma whose source label is \texttt{lem:res2} in the article (source0.tex lines 554--575, source label \texttt{lem:res2}), delivered together with the article's own complete original proof of it (source0.tex lines 578--622, one \texttt{proof} environment of 45 lines). No mathematics is written, repaired or re-derived: statement and proof are the article's own text line for line. Required context retained uncounted: eq:Ak (lines 472--481); lem:res1 (lines 494--510). Every definition, display and label of those lines is retained unchanged. Omitted from this package and not redistributed: 1--471, 482--493, 511--553, 576--577, 623--1264. Nothing inside a retained excerpt is omitted or altered: every retained body is delivered exactly as the article's lines are written, so no omission receipt of any kind accompanies this package. Citations: the retained excerpts carry no citation command at all, so no bibliography entry of the article is transcribed and no reference is redistributed. Reference adaptation: none was needed and none was made; every reference inside the delivered unit resolves inside it, so no unresolved reference or citation is produced. Printed slips kept literal: no printed word, symbol, hypothesis or number of the counted statement or of its proof was altered. Layout and class: the article's own class is \texttt{amsart} with packages including geometry, fontenc, babel, amsthm, amsfonts, amsmath, amssymb, multicol, mathtools, enumerate, hyperref, array, enumitem, color; the wrapper is the lane's pinned \texttt{amsart} editorial wrapper at 10pt on A4, which restates the theorem-like environments the retained text needs and adds the article's own macro definitions that the retained text needs (\texttt{\textbackslash NN}, \texttt{\textbackslash Q}, \texttt{\textbackslash Res}, \texttt{\textbackslash geq}, \texttt{\textbackslash leq}), with extra wrapper packages (bm, bbm, dsfont, xspace, cancel, mathtools). The delivered numbering is the unit's own. Assets: no retained span contains a figure environment, a graphics-inclusion command, a PDF-page inclusion, a table or a TikZ picture; no image file is copied into this package, the package declares no asset dependency and no image file is redistributed with it. Licence: version arXiv:2411.09264v2 of this article is distributed under the Creative Commons Attribution 4.0 International licence, as recorded in the arXiv licence metadata for this exact version and shown on the abstract page https://arxiv.org/abs/2411.09264v2. A full rights notice is delivered at licenses/arxiv-2411.09264-CC-BY-4.0.txt. Characters: every retained excerpt is pure ASCII, so the delivered engine, XeLaTeX with its default Latin Modern text font, reports no missing character.
\begin{equation}\label{eq:Ak}
A_k=
\begin{pmatrix}
 a& 0& \cdots  & 0 \\
b& a&  \cdots& 0     \\
c& b& \cdots  & 0\\
0& c& \cdots  & a \\
 \vdots&\vdots & \ddots  & b  \\
  0&  0&\cdots & c
\end{pmatrix}.\end{equation}

\begin{lemma}\label{lem:res1}
Let $1\leq i<j\leq n+1$. Then there exist $X,Y\in\Q[r_0,\dots, r_n,a,b,c]$ depending on the choice of $i,j$ such that 
$$
\Res(R,Q)=
\det
\begin{pmatrix}
\vdots & \vdots & &&&& \\
X     & 0     & &&&&\\
0     & X     & &&&&\\
\vdots & \vdots &&&A_n&& \\
Y     & 0     &  &&&&\\ 
0     & Y     &     &&&&\\
\vdots & \vdots & &&&&
\end{pmatrix},
$$
where $X$ appears in the $(i,1)$-entry and $(i+1,2)$-entry of the matrix,  and $Y$ appears in the $(j,1)$-entry and $(j+1,2)$-entry.
\end{lemma}

\begin{lemma}\label{lem:res2}
Let $1\leq i<j\leq n+1$.
Then there exist $X,Y\in\Q[r_0,\dots, r_n,a,b,c]$ depending on the choice of $i,j$ such that 
\[
\Res(R,Q)
=a^{i-1}c^{n-j+1}\left(c^{j-i}X^2+(-1)^{j-i}(d_{j-i}-acd_{j-i-2})XY+a^{j-i}Y^2\right),
\]
where 
$d_{-1}=0$,  $d_0=1$, and $d_k$ is the $k\times k$ determinant 
\[d_k=\det
\begin{pmatrix}
b & a & 0& \cdots  & &  & \\
c & b& a&  \cdots & &  0&\\
0 & c& b& \cdots &  &  & \\
\vdots & \vdots &\vdots & \ddots  &\vdots &\vdots &\vdots \\ 
   &    &  & \dots& b& a&  0\\
   & 0  &   &\dots& c& b&  a\\
   &    &  &\dots & 0& c&  b
\end{pmatrix},
\]
for $k\in \NN$.
\end{lemma}

\begin{proof}
Recall the definition \eqref{eq:Ak} of $A_k$, 
and the subsequent definition of $B_{k,i,j}$.
(In particular $d_k=\det B_{k,1,k+2}$ in this notation, for $k\geq 1$.)
 Let $i,j$ be fixed such that  $1\leq i<j\leq n+1$. Then 
it follows from Lemma \ref{lem:res1} that 
\[
\Res(R,Q)=(-1)^{i+1}X\det M_{i,1}+(-1)^{j+1}Y\det M_{j,1},
\]
where $M_{k,1}$ is the $(n+1)\times(n+1)$ matrix by removing the $k$-th row and the first column.
Now the matrix $M_{i,1}$ has an $X$ in its $(i,1)$-entry and a $Y$ in its $(j,1)$-entry, so
\[\det M_{i,1}=(-1)^{i+1}X\det B_{n,i,i+1}+(-1)^{j+1}Y\det B_{n,i,j+1}.\]
Similarly, 
the matrix $M_{j,1}$ has an $X$ in its $(i+1,1)$-entry if $j>i+1$, and a $Y$ in its $(j,1)$-entry, so
\[\det M_{j,1}=(-1)^{i+2}X\mathbf{1}_{j> i+1}\det B_{n,i+1,j}+(-1)^{j+1} Y\det B_{n,j,j+1}.\]
Putting everything together, we obtain
\begin{equation}\label{eq:Res1}
\begin{split}
\Res(R,Q)
=~&
X^2\det B_{n,i,i+1}
+Y^2 \det B_{n,j,j+1}\\
&+(-1)^{j-i}\left(\det B_{n,i,j+1}-\mathbf{1}_{j> i+1}\det B_{n,i+1,j}\right)XY.
\end{split}
\end{equation}


It remains to compute $\det B_{n,i,j+1}$ for 
 integers $1\leq i\leq j\leq n+1$. 
If $i>1$,  Laplace expansion along the first row of $B_{n,i,j+1}$ gives
\[ \det B_{n,i,j+1}=a\det B_{n-1,i-1,j}.\]
If $j\leq n$, Laplace expansion along the final row of $B_{n,i,j+1}$ gives
\[ \det B_{n,i,j+1}=c\det B_{n-1,i,j+1}.\]
Applying these relations recursively, we easily obtain
\[ \det B_{n,i,j+1}
=a^{i-1}\det B_{n-i+1,1,j-i+2}
=a^{i-1}c^{n-j+1}
\det B_{j-i,1,j-i+2}.
\]
Recalling the definition of $d_k$ from the statement, 
we see that 
\[\det B_{n,i,j+1}=a^{i-1}c^{n-j+1}d_{j-i},
\]
which  easily leads to the statement of the lemma, on insertion into  \eqref{eq:Res1}.
\end{proof}

\end{document}
